\documentclass{article}
\usepackage[T1]{fontenc}
\usepackage{lmodern}
\usepackage{graphicx}
\usepackage{amsmath,amssymb}
\usepackage[a4paper,margin=2cm]{geometry}
\usepackage[absolute,overlay]{textpos}
\setlength{\TPHorizModule}{1mm}
\setlength{\TPVertModule}{1mm}
\usepackage[indonesian]{babel}
\usepackage{hyperref}
\setlength{\emergencystretch}{2em}
% unit: Halaman 60

\begin{document}

% === Halaman 60 (python/native) ===

2.  Algoritma kriptografi nir-simetri (asymmetric-key cryptography)

•  Kunci enkripsi  kunci dekripsi (K1  K2)

     Kunci enkripsi → tidak rahasia (public key)

     Kunci dekripsi → rahasia (private key)

• Mulai ditemukan sejak tahun 1976

60

Kunci publik, K1

Enkripsi

EK1 (P) = C

Dekripsi

DK2 (C) = P

Cipherteks, C

Plainteks, P

Plainteks, P

Kunci privat, K2

60

-

RSA (Rivest-Shamir-Adleman)

-

ElGamal

-

DSA

-

Diffie-Hellman 

-

Mercke Knapsack Algorithm

-

Rabin

-

EPOC

-

Mc Eliece

-

XTR

-

ECC (Elliptic Curve Cryptography)

Nama lain: Kriptografi kunci –publik 

(public-key cryptography)

\end{document}
